\chapter{Introduction}

In this, hopefully, short chat I will introduce you to GitHub. There will be a number of separate methods of using GitHub, as follows:

\begin{itemize}
    \item \emph{Very simple}: downloading a ``release'' with no interest in the source code.
    \item \emph{Simple}: grabbing someone's source code, as a Zip file, for our own use. This might not include a runnable binary---download a release as well if necessary.
    \item \emph{Medium}: grabbing someone's source code with the possibility of updating it, but \emph{only} for our own use---we will ``clone'' a repository.
    \item \emph{Slightly advanced}: grabbing someone's code with the intention of updating it for ourselves and everyone else to use---we will ``fork'' a repository, then ``clone'' it.
    \item \emph{Advanced}: here we look at creating websites for our \emph{organisation}\footnote{By \emph{organisation} I mean user or group as setup by ourselves, see later.} and also for each individual project we have in GitHub.
\end{itemize}

\section{What is Git}
Git is a distributed version control system, created by Linus Torvalds the original creator of Linux.

The code for the Linux Kernel was held in a proprietary version control system but was given special permissions to use it for no, or little, charge.

Once fully embedded, the company decided to do the dirty on the Linux Kernel folks, so Linus went ``stuff you'' and wrote the original Git in a weekend.

Since then, Git has become probably the worlds best known and most used [distributed] version control system---there are a couple of others. \emph{Mercury} for example.

A distributed version control system means that anyone using it for a project has a full copy of the source code \emph{plus} the entire history of changes, releases and such like.

Work can be done off-line and away from the office, and uploaded later when back in contact, without suffering from \emph{conflicts} if someone else edited the same files as you have.

Obviously, conflicts will arise if two or more folk edit the same \emph{line(s)} of code. Git will help resolve the mess if it occurs.


\section{What is GitHub}
GitHub, \url{https://github.com}, is an online code repository, now owned by Microsoft, which offers individuals and/or organisations:

\begin{itemize}
    \item Distributed version control of documents, code and almost anything else. Off-site backups in other words, plus the ability to undo horrendous cock-ups!
    \item Public and private repositories for projects.
    \item Web Pages for the organisation or individual.
    \item Web pages for individual projects.
    \item And much more: Wikis, Gists, Notifications, team working \ldots
\end{itemize}

I won't be covering everything GitHub offers tonight, some of us will have beds to go to.

Other \emph{Git} style code repositories are also available, the ones I'm familiar with being:

\begin{itemize}
    \item GitLab, \url{https://gitlab.com}, is very similar to GitHub and not owned by Microsoft. My books live here. GitLab offers both private and public repositories.
    \item CodeBerg - A recent arrival for people who don't want Microsoft anywhere near their code especially with their so-called ``AI'' which literally steals anything and everything and claims it for its own!

    CodeBerg, \url{https://codeberg.org}, only offers public repositories as it caters to free and Open Source projects only.
\end{itemize}